\section{Introduction}

The transition to post-quantum key agreement is underway. NIST has
standardized ML-KEM, other state-level bodies are providing their own
cryptographic recommendations, and organizations are beginning to update to
to protect against future quantum adversaries. However, these algorithms are
relatively new, and there remains a possibility of unexpected
cryptanalytic breakthroughs or vulnerable implementations~\cite{EC:CasDec23}. A natural hedge
against this risk is to deploy hybrid constructions that combine post-quantum
and classical algorithms, maintaining security as long as at least one
component remains secure.

This ``defense in depth'' approach is appealing: if the post-quantum KEM is
broken by classical cryptanalysis or a software bug, the classical component
provides classical security; if the classical component is broken by a
quantum computer, the post-quantum component provides quantum
resistance. However, constructing provably secure hybrid KEMs requires
care. The combiner function must preserve security properties under the right
assumptions, and different designs involve different trade-offs between the
cryptographic assumptions of the components and efficiency of the whole
scheme.

The Internet Research Task Force (IRTF)'s Crypto Forum Research Group (CFRG) is working towards finalizing a recommendation for specific constructions of hybrid KEMs. The current document~\cite{irtf-cfrg-hybrid-kems-05} gives four constructions named CG, CK, UG, and UK. The 

\paragraph{Our Contributions.} In this paper, we introduce and analyze $\UG$
(Universal, Group), a generic hybrid KEM framework that combines a
post-quantum KEM with a classically-secure nominal group. Unlike other
constructions, we rely directly on the Strong Diffie-Hellman (SDH) security
of the group, and include all public encapsulation keys and ciphertexts, so
that we only rely on the IND-CCA security of the PQ KEM. Our main
contributions are:

We present a hybrid KEM construction that concatenates shared secrets,
ciphertexts, and encapsulation keys from both components, then applies a key
derivation function (KDF) to produce the final shared secret. The
construction is deliberately conservative: by including all public values in
the KDF input, $\UG$ makes no special assumptions about its component schemes
beyond IND-CCA security of the post-quantum KEM and SDH security of the
nominal group to achieve IND-CCA security of the hybrid KEM, when either of
the component schemes lose their security.

We provide two complementary proofs that $\UG$ is IND-CCA secure. In the
pre-quantum case, we prove security in the Random Oracle Model: if the
nominal group is SDH secure against classical adversaries, then so is
$\UG$. This holds even if the post-quantum KEM is broken (without
compromising the nominal group's secrets). We model the KDF as a Random
Oracle. In the post-quantum case, we show security in the Standard Model: if
the post-quantum KEM is IND-CCA secure against quantum adversaries, then so
is $\UG$. We model the KDF as a pseudorandom function, keyed by the
post-quantum shared secret, so breaking $\UG$ requires breaking the
post-quantum KEM. This holds even if the nominal group SDH security is
completely broken.

Notably, $\UG$ does not require ciphertext second-preimage resistance
(C2PRI)\cite{CiC:BCDKSV24} from its components— SDH and IND-CCA security
suffices, nor that the classical component be an IND-CCA KEM. Thus, $\UG$
can be instantiated with diverse combinations of elliptic curves as nominal
groups, and any IND-CCA secure post-quantum KEM.

Our work complements recent hybrid KEM designs that optimize for performance
by leveraging additional assumptions. $\UG$ represents a conservative point in
the design space: it achieves strong security guarantees with minimal
assumptions, at the cost of hashing additional public data, while eliminating
layers of abstraction.

\paragraph{Structure of this paper.}
Notation and definitions are introduced \autoref{sec:prelims}.  The $\UG$
framework is introduced in \autoref{sec:ug}; our generic construction --
named Heavy Quantum Bomber -- is a general framework that can be instantiated
with different component algorithms.  In \autoref{sec:ug-proofs} we prove
the IND-CCA security when the different components fail in different ways.
